\documentclass[
	ngerman,
	toc=listof, % Abbildungsverzeichnis sowie Tabellenverzeichnis in das Inhaltsverzeichnis aufnehmen
	toc=bibliography, % Literaturverzeichnis in das Inhaltsverzeichnis aufnehmen
	footnotes=multiple, % Trennen von direkt aufeinander folgenden Fußnoten
	parskip=half, % vertikalen Abstand zwischen Absätzen verwenden anstatt horizontale Einrückung von Folgeabsätzen
	numbers=noendperiod % Den letzten Punkt nach einer Nummerierung entfernen (nach DIN 5008)
]{scrartcl}

\usepackage[utf8]{inputenc} % muss als erstes eingebunden werden, da Meta/Packages ggfs. Sonderzeichen enthalten

% !TEX root = Projektdokumentation.tex

\newcommand{\titel}{Webapplikation zur Verwaltung der Zugriffsrechte der Benutzeraccounts externer Services}
\newcommand{\untertitel}{Webbasiertes Tool zur Unterstützung der Entwickler}
\newcommand{\kompletterTitel}{\titel{} -- \untertitel}

\newcommand{\autorName}{Farah Schüller}
\newcommand{\autorAnschrift}{Straße 123}
\newcommand{\autorOrt}{1337 Stadt}

\newcommand{\betriebLogo}{LogoBetrieb.pdf}
\newcommand{\betriebName}{\textsc{SUSE Linux} GmbH}
\newcommand{\betriebAnschrift}{Maxfeldstraße 5}
\newcommand{\betriebOrt}{90409 Nürnberg}

\newcommand{\ausbildungsberuf}{Fachinformatiker für Anwendungsentwicklung}
\newcommand{\betreff}{Dokumentation zur betrieblichen Projektarbeit}
\newcommand{\pruefungstermin}{Winter 2016}
\newcommand{\abgabeOrt}{Nürnberg}
\newcommand{\abgabeTermin}{15.12.2016}
 % Metadaten zu diesem Dokument (Autor usw.)
% !TEX root = ../Projektdokumentation.tex

% Anpassung an Landessprache ---------------------------------------------------
\usepackage{babel}

% Umlaute ----------------------------------------------------------------------
%   Umlaute/Sonderzeichen wie äüöß direkt im Quelltext verwenden (CodePage).
%   Erlaubt automatische Trennung von Worten mit Umlauten.
% ------------------------------------------------------------------------------
\usepackage[T1]{fontenc}
\usepackage{textcomp} % Euro-Zeichen etc.

% Schrift ----------------------------------------------------------------------
\usepackage{lmodern} % bessere Fonts
%\usepackage{relsize} % Schriftgröße relativ festlegen

% Tabellen ---------------------------------------------------------------------
\PassOptionsToPackage{table}{xcolor}
\usepackage{tabularx}
% für lange Tabellen
\usepackage{longtable}
\usepackage{array}
\usepackage{ragged2e}
\usepackage{lscape}
\newcolumntype{w}[1]{>{\raggedleft\hspace{0pt}}p{#1}} % Spaltendefinition rechtsbündig mit definierter Breite

% Grafiken ---------------------------------------------------------------------
\usepackage[dvips,final]{graphicx} % Einbinden von JPG-Grafiken ermöglichen
\usepackage{graphics} % keepaspectratio
\usepackage{floatflt} % zum Umfließen von Bildern
\graphicspath{{Bilder/}} % hier liegen die Bilder des Dokuments

% Sonstiges --------------------------------------------------------------------
\usepackage[titles]{tocloft} % Inhaltsverzeichnis DIN 5008 gerecht einrücken
\usepackage{amsmath,amsfonts} % Befehle aus AMSTeX für mathematische Symbole
\usepackage{enumitem} % anpassbare Enumerates/Itemizes
\usepackage{xspace} % sorgt dafür, dass Leerzeichen hinter parameterlosen Makros nicht als Makroendezeichen interpretiert werden

\usepackage{makeidx} % für Index-Ausgabe mit \printindex
\usepackage[printonlyused]{acronym} % es werden nur benutzte Definitionen aufgelistet

% Einfache Definition der Zeilenabstände und Seitenränder etc.
\usepackage{setspace}
\usepackage{geometry}
\usepackage{lastpage}
% \usepackage{pageslts}

% Symbolverzeichnis
\usepackage[intoc]{nomencl}
\let\abbrev\nomenclature
\renewcommand{\nomname}{Abkürzungsverzeichnis}
\setlength{\nomlabelwidth}{.25\hsize}
\renewcommand{\nomlabel}[1]{#1 \dotfill}
\setlength{\nomitemsep}{-\parsep}

\usepackage{varioref} % Elegantere Verweise. „auf der nächsten Seite“
\usepackage{url} % URL verlinken, lange URLs umbrechen etc.

\usepackage{chngcntr} % fortlaufendes Durchnummerieren der Fußnoten
% \usepackage[perpage]{footmisc} % Alternative: Nummerierung der Fußnoten auf jeder Seite neu

\usepackage{ifthen} % bei der Definition eigener Befehle benötigt
\usepackage{todonotes} % definiert u.a. die Befehle \todo und \listoftodos
\usepackage[square]{natbib} % wichtig für korrekte Zitierweise

% PDF-Optionen -----------------------------------------------------------------
\usepackage{pdfpages}
\pdfminorversion=5 % erlaubt das Einfügen von pdf-Dateien bis Version 1.7, ohne eine Fehlermeldung zu werfen (keine Garantie für fehlerfreies Einbetten!)
\usepackage[
    bookmarks,
    bookmarksnumbered,
    bookmarksopen=true,
    bookmarksopenlevel=1,
    colorlinks=true,
% diese Farbdefinitionen zeichnen Links im PDF farblich aus
    linkcolor=AOBlau, % einfache interne Verknüpfungen
    anchorcolor=AOBlau,% Ankertext
    citecolor=AOBlau, % Verweise auf Literaturverzeichniseinträge im Text
    filecolor=AOBlau, % Verknüpfungen, die lokale Dateien öffnen
    menucolor=AOBlau, % Acrobat-Menüpunkte
    urlcolor=AOBlau,
% diese Farbdefinitionen sollten für den Druck verwendet werden (alles schwarz)
    %linkcolor=black, % einfache interne Verknüpfungen
    %anchorcolor=black, % Ankertext
    %citecolor=black, % Verweise auf Literaturverzeichniseinträge im Text
    %filecolor=black, % Verknüpfungen, die lokale Dateien öffnen
    %menucolor=black, % Acrobat-Menüpunkte
    %urlcolor=black,
%
    %backref, % Quellen werden zurück auf ihre Zitate verlinkt
    pdftex,
    plainpages=false, % zur korrekten Erstellung der Bookmarks
    pdfpagelabels=true, % zur korrekten Erstellung der Bookmarks
    hypertexnames=false, % zur korrekten Erstellung der Bookmarks
    linktocpage % Seitenzahlen anstatt Text im Inhaltsverzeichnis verlinken
]{hyperref}
% Befehle, die Umlaute ausgeben, führen zu Fehlern, wenn sie hyperref als Optionen übergeben werden
\hypersetup{
    pdftitle={\titel -- \untertitel},
    pdfauthor={\autorName},
    pdfcreator={\autorName},
    pdfsubject={\titel -- \untertitel},
    pdfkeywords={\titel -- \untertitel},
}


% zum Einbinden von Programmcode -----------------------------------------------
\usepackage{listings}
\usepackage{xcolor}
\definecolor{hellgelb}{rgb}{1,1,0.9}
\definecolor{colKeys}{rgb}{0,0,1}
\definecolor{colIdentifier}{rgb}{0,0,0}
\definecolor{colComments}{rgb}{0,0.5,0}
\definecolor{colString}{rgb}{1,0,0}
\lstset{
    float=hbp,
	basicstyle=\footnotesize,
    identifierstyle=\color{colIdentifier},
    keywordstyle=\color{colKeys},
    stringstyle=\color{colString},
    commentstyle=\color{colComments},
    backgroundcolor=\color{hellgelb},
    columns=flexible,
    tabsize=2,
    frame=single,
    extendedchars=true,
    showspaces=false,
    showstringspaces=false,
    numbers=left,
    numberstyle=\tiny,
    breaklines=true,
    breakautoindent=true,
	captionpos=b,
}
\lstdefinelanguage{cs}{
	sensitive=false,
	morecomment=[l]{//},
	morecomment=[s]{/*}{*/},
	morestring=[b]",
	morekeywords={
		abstract,event,new,struct,as,explicit,null,switch
		base,extern,object,this,bool,false,operator,throw,
		break,finally,out,true,byte,fixed,override,try,
		case,float,params,typeof,catch,for,private,uint,
		char,foreach,protected,ulong,checked,goto,public,unchecked,
		class,if,readonly,unsafe,const,implicit,ref,ushort,
		continue,in,return,using,decimal,int,sbyte,virtual,
		default,interface,sealed,volatile,delegate,internal,short,void,
		do,is,sizeof,while,double,lock,stackalloc,
		else,long,static,enum,namespace,string},
}
\lstdefinelanguage{natural}{
	sensitive=false,
	morecomment=[l]{/*},
	morestring=[b]",
	morestring=[b]',
	alsodigit={-,*},
	morekeywords={
		DEFINE,DATA,LOCAL,END-DEFINE,WRITE,CALLNAT,PARAMETER,USING,
		IF,NOT,END-IF,ON,*ERROR-NR,ERROR,END-ERROR,ESCAPE,ROUTINE,
		PERFORM,SUBROUTINE,END-SUBROUTINE,CONST,END-FOR,END,FOR,RESIZE,
		ARRAY,TO,BY,VALUE,RESET,COMPRESS,INTO,EQ},
}
\lstdefinelanguage{php}{
	sensitive=false,
	morecomment=[l]{/*},
	morestring=[b]",
	morestring=[b]',
	alsodigit={-,*},
	morekeywords={
		abstract,and,array,as,break,case,catch,cfunction,class,clone,const,
		continue,declare,default,do,else,elseif,enddeclare,endfor,endforeach,
		endif,endswitch,endwhile,extends,final,for,foreach,function,global,
		goto,if,implements,interface,instanceof,namespace,new,old_function,or,
		private,protected,public,static,switch,throw,try,use,var,while,xor
		die,echo,empty,exit,eval,include,include_once,isset,list,require,
		require_once,return,print,unset},
}
 % verwendete Packages
% !TEX root = ../Projektdokumentation.tex

% Seitenränder -----------------------------------------------------------------
\setlength{\topskip}{\ht\strutbox} % behebt Warnung von geometry
\geometry{a4paper,left=20mm,right=20mm,top=25mm,bottom=35mm}

\usepackage[
	automark, % Kapitelangaben in Kopfzeile automatisch erstellen
	headsepline, % Trennlinie unter Kopfzeile
	ilines % Trennlinie linksbündig ausrichten
]{scrpage2}

% Kopf- und Fußzeilen ----------------------------------------------------------
\pagestyle{scrheadings}
% chapterpagestyle gibt es nicht in scrartcl
%\renewcommand{\chapterpagestyle}{scrheadings}
\clearscrheadfoot

% Kopfzeile
\renewcommand{\headfont}{\normalfont} % Schriftform der Kopfzeile
\ihead{\large{\textsc{\titel}}\\ \small{\untertitel} \\[2ex] \textit{\headmark}}
\chead{}
\ohead{\includegraphics[scale=0.15]{\betriebLogo}}
\setlength{\headheight}{25mm} % Höhe der Kopfzeile
%\setheadwidth[0pt]{textwithmarginpar} % Kopfzeile über den Text hinaus verbreitern (falls Logo den Text überdeckt)

% Fußzeile
\ifoot{\autorName}
\cfoot{}
\ofoot{\pagemark}

% Überschriften nach DIN 5008 in einer Fluchtlinie
% ------------------------------------------------------------------------------

% Abstand zwischen Nummerierung und Überschrift definieren
% > Schön wäre hier die dynamische Berechnung des Abstandes in Abhängigkeit
% > der Verschachtelungstiefe des Inhaltsverzeichnisses
\newcommand{\headingSpace}{1.5cm}

% Abschnittsüberschriften im selben Stil wie beim Inhaltsverzeichnis einrücken
\renewcommand*{\othersectionlevelsformat}[3]{
  \makebox[\headingSpace][l]{#3\autodot}
}

% Für die Einrückung wird das Paket tocloft benötigt
%\cftsetindents{chapter}{0.0cm}{\headingSpace}
\cftsetindents{section}{0.0cm}{\headingSpace}
\cftsetindents{subsection}{0.0cm}{\headingSpace}
\cftsetindents{subsubsection}{0.0cm}{\headingSpace}
\cftsetindents{figure}{0.0cm}{\headingSpace}
\cftsetindents{table}{0.0cm}{\headingSpace}


% Allgemeines
% ------------------------------------------------------------------------------

\onehalfspacing % Zeilenabstand 1,5 Zeilen
\frenchspacing % erzeugt ein wenig mehr Platz hinter einem Punkt

% Schusterjungen und Hurenkinder vermeiden
\clubpenalty = 10000
\widowpenalty = 10000
\displaywidowpenalty = 10000

% Quellcode-Ausgabe formatieren
\lstset{numbers=left, numberstyle=\tiny, numbersep=5pt, breaklines=true}
\lstset{emph={square}, emphstyle=\color{red}, emph={[2]root,base}, emphstyle={[2]\color{blue}}}

\counterwithout{footnote}{section} % Fußnoten fortlaufend durchnummerieren
\setcounter{tocdepth}{\subsubsectionlevel} % im Inhaltsverzeichnis werden die Kapitel bis zum Level der subsubsection übernommen
\setcounter{secnumdepth}{\subsubsectionlevel} % Kapitel bis zum Level der subsubsection werden nummeriert

% Aufzählungen anpassen
\renewcommand{\labelenumi}{\arabic{enumi}.}
\renewcommand{\labelenumii}{\arabic{enumi}.\arabic{enumii}.}
\renewcommand{\labelenumiii}{\arabic{enumi}.\arabic{enumii}.\arabic{enumiii}}

% Tabellenfärbung:
\definecolor{heading}{rgb}{0.64,0.78,0.86}
\definecolor{odd}{rgb}{0.9,0.9,0.9}
 % Definitionen zum Aussehen der Seiten
% !TEX root = ../Projektdokumentation.tex

% Abkürzungen, ggfs. mit korrektem Leerraum
\newcommand{\bs}{$\backslash$\xspace}
\newcommand{\bspw}{bspw.\xspace}
\newcommand{\bzw}{bzw.\xspace}
\newcommand{\ca}{ca.\xspace}
\newcommand{\dahe}{\mbox{d.\,h.}\xspace}
\newcommand{\etc}{etc.\xspace}
\newcommand{\eur}[1]{\mbox{#1\,\texteuro}\xspace}
\newcommand{\evtl}{evtl.\xspace}
\newcommand{\ggfs}{ggfs.\xspace}
\newcommand{\Ggfs}{Ggfs.\xspace}
\newcommand{\gqq}[1]{\glqq{}#1\grqq{}}
\newcommand{\inkl}{inkl.\xspace}
\newcommand{\insb}{insb.\xspace}
\newcommand{\ua}{\mbox{u.\,a.}\xspace}
\newcommand{\usw}{usw.\xspace}
\newcommand{\Vgl}{vgl.\xspace}
\newcommand{\zB}{\mbox{z.\,B.}\xspace}
\newcommand{\anf}[1]{\glqq#1\grqq}

% Befehle für häufig anfallende Aufgaben
\newcommand{\Abbildung}[1]{\autoref{fig:#1}}
\newcommand{\Anhang}[1]{\appendixname{}~\ref{#1}: \nameref{#1} \vpageref{#1}}
\newcommand{\includegraphicsKeepAspectRatio}[2]{\includegraphics[width=#2\textwidth,height=#2\textheight,keepaspectratio]{#1}}
\newcommand{\Zitat}[2][\empty]{\ifthenelse{\equal{#1}{\empty}}{\citep{#2}}{\citep[#1]{#2}}}
\newcommand{\Autor}[1]{\textsc{#1}} % zum Ausgeben von Autoren
\newcommand{\itemd}[2]{\item{\textbf{#1}}\\{#2}} % erzeugt ein Listenelement mit fetter Überschrift

% füge Code-Snippet mit Caption ein
\newcommand{\snip}[3]
{
\lstinputlisting[caption=#1, firstline=#2, lastline=#3]{Listings/#1}
}


% füge Abbildungen mit Caption ein
\newcommand{\abb}[4]
{\begin{figure}[htb]
\centering
\singlespacing
\includegraphicsKeepAspectRatio{Bilder/#3}{#4}
\caption{#1}
\label{#2}
\end{figure}}

% fügt Tabellen aus einer TEX-Datei ein
\newcommand{\tabelle}[3] % Parameter: caption, label, file
{\begin{table}[htbp]
\centering
\singlespacing
\input{Tabellen/#3}
\caption{#1}
\label{#2}
\end{table}}

\newcommand{\tabelleAnhang}[1] % Parameter: file
{\begin{center}
\singlespacing
\input{Tabellen/#1}
\end{center}}

% einfaches Wechseln der Schrift, z.B.: \changefont{cmss}{sbc}{n}
\newcommand{\changefont}[3]{\fontfamily{#1} \fontseries{#2} \fontshape{#3} \selectfont}

% Verwendung analog zu \includegraphics
\newlength{\myx} % Variable zum Speichern der Bildbreite
\newlength{\myy} % Variable zum Speichern der Bildhöhe
\newcommand\includegraphicstotab[2][\relax]{%
% Abspeichern der Bildabmessungen
\settowidth{\myx}{\includegraphics[{#1}]{#2}}%
\settoheight{\myy}{\includegraphics[{#1}]{#2}}%
% das eigentliche Einfügen
\parbox[c][1.1\myy][c]{\myx}{%
\includegraphics[{#1}]{#2}}%
}

\definecolor{AOBlau}{rgb}{0, 0, 0}

% verschiedene Befehle um Wörter semantisch auszuzeichnen ----------------------
\newcommand{\Index}[2][\empty]{\ifthenelse{\equal{#1}{\empty}}{\index{#2}#2}{\index{#1}#2}}
\newcommand{\Fachbegriff}[2][\empty]{\ifthenelse{\equal{#1}{\empty}}{\textit{\Index{#2}}}{\textit{\Index[#1]{#2}}}}
\newcommand{\NeuerBegriff}[2][\empty]{\ifthenelse{\equal{#1}{\empty}}{\textbf{\Index{#2}}}{\textbf{\Index[#1]{#2}}}}

\newcommand{\Ausgabe}[1]{\texttt{#1}}
\newcommand{\Eingabe}[1]{\texttt{#1}}
\newcommand{\Code}[1]{\texttt{#1}}
\newcommand{\Datei}[1]{\texttt{#1}}

\newcommand{\Assembly}[1]{\textsf{#1}}
\newcommand{\Klasse}[1]{\textsf{#1}}
\newcommand{\Methode}[1]{\textsf{#1}}
\newcommand{\Attribut}[1]{\textsf{#1}}

\newcommand{\Datentyp}[1]{\textsf{#1}}
\newcommand{\XMLElement}[1]{\textsf{#1}}
\newcommand{\Webservice}[1]{\textsf{#1}}

\newcommand{\Refactoring}[1]{\Fachbegriff{#1}}
\newcommand{\CodeSmell}[1]{\Fachbegriff{#1}}
\newcommand{\Metrik}[1]{\Fachbegriff{#1}}
\newcommand{\DesignPattern}[1]{\Fachbegriff{#1}}
 % eigene allgemeine Befehle, die z.B. die Arbeit mit LaTeX erleichtern
% % !TEX root = ../Projektdokumentation.tex

% Abkürzungen, ggfs. mit korrektem Leerraum
\newcommand{\bs}{$\backslash$\xspace}
\newcommand{\bspw}{bspw.\xspace}
\newcommand{\bzw}{bzw.\xspace}
\newcommand{\ca}{ca.\xspace}
\newcommand{\dahe}{\mbox{d.\,h.}\xspace}
\newcommand{\etc}{etc.\xspace}
\newcommand{\eur}[1]{\mbox{#1\,\texteuro}\xspace}
\newcommand{\evtl}{evtl.\xspace}
\newcommand{\ggfs}{ggfs.\xspace}
\newcommand{\Ggfs}{Ggfs.\xspace}
\newcommand{\gqq}[1]{\glqq{}#1\grqq{}}
\newcommand{\inkl}{inkl.\xspace}
\newcommand{\insb}{insb.\xspace}
\newcommand{\ua}{\mbox{u.\,a.}\xspace}
\newcommand{\usw}{usw.\xspace}
\newcommand{\Vgl}{vgl.\xspace}
\newcommand{\zB}{\mbox{z.\,B.}\xspace}
\newcommand{\anf}[1]{\glqq#1\grqq}

% Befehle für häufig anfallende Aufgaben
\newcommand{\Abbildung}[1]{\autoref{fig:#1}}
\newcommand{\Anhang}[1]{\appendixname{}~\ref{#1}: \nameref{#1} \vpageref{#1}}
\newcommand{\includegraphicsKeepAspectRatio}[2]{\includegraphics[width=#2\textwidth,height=#2\textheight,keepaspectratio]{#1}}
\newcommand{\Zitat}[2][\empty]{\ifthenelse{\equal{#1}{\empty}}{\citep{#2}}{\citep[#1]{#2}}}
\newcommand{\Autor}[1]{\textsc{#1}} % zum Ausgeben von Autoren
\newcommand{\itemd}[2]{\item{\textbf{#1}}\\{#2}} % erzeugt ein Listenelement mit fetter Überschrift

% füge Code-Snippet mit Caption ein
\newcommand{\snip}[3]
{
\lstinputlisting[caption=#1, firstline=#2, lastline=#3]{Listings/#1}
}


% füge Abbildungen mit Caption ein
\newcommand{\abb}[4]
{\begin{figure}[htb]
\centering
\singlespacing
\includegraphicsKeepAspectRatio{Bilder/#3}{#4}
\caption{#1}
\label{#2}
\end{figure}}

% fügt Tabellen aus einer TEX-Datei ein
\newcommand{\tabelle}[3] % Parameter: caption, label, file
{\begin{table}[htbp]
\centering
\singlespacing
\input{Tabellen/#3}
\caption{#1}
\label{#2}
\end{table}}

\newcommand{\tabelleAnhang}[1] % Parameter: file
{\begin{center}
\singlespacing
\input{Tabellen/#1}
\end{center}}

% einfaches Wechseln der Schrift, z.B.: \changefont{cmss}{sbc}{n}
\newcommand{\changefont}[3]{\fontfamily{#1} \fontseries{#2} \fontshape{#3} \selectfont}

% Verwendung analog zu \includegraphics
\newlength{\myx} % Variable zum Speichern der Bildbreite
\newlength{\myy} % Variable zum Speichern der Bildhöhe
\newcommand\includegraphicstotab[2][\relax]{%
% Abspeichern der Bildabmessungen
\settowidth{\myx}{\includegraphics[{#1}]{#2}}%
\settoheight{\myy}{\includegraphics[{#1}]{#2}}%
% das eigentliche Einfügen
\parbox[c][1.1\myy][c]{\myx}{%
\includegraphics[{#1}]{#2}}%
}

\definecolor{AOBlau}{rgb}{0, 0, 0}

% verschiedene Befehle um Wörter semantisch auszuzeichnen ----------------------
\newcommand{\Index}[2][\empty]{\ifthenelse{\equal{#1}{\empty}}{\index{#2}#2}{\index{#1}#2}}
\newcommand{\Fachbegriff}[2][\empty]{\ifthenelse{\equal{#1}{\empty}}{\textit{\Index{#2}}}{\textit{\Index[#1]{#2}}}}
\newcommand{\NeuerBegriff}[2][\empty]{\ifthenelse{\equal{#1}{\empty}}{\textbf{\Index{#2}}}{\textbf{\Index[#1]{#2}}}}

\newcommand{\Ausgabe}[1]{\texttt{#1}}
\newcommand{\Eingabe}[1]{\texttt{#1}}
\newcommand{\Code}[1]{\texttt{#1}}
\newcommand{\Datei}[1]{\texttt{#1}}

\newcommand{\Assembly}[1]{\textsf{#1}}
\newcommand{\Klasse}[1]{\textsf{#1}}
\newcommand{\Methode}[1]{\textsf{#1}}
\newcommand{\Attribut}[1]{\textsf{#1}}

\newcommand{\Datentyp}[1]{\textsf{#1}}
\newcommand{\XMLElement}[1]{\textsf{#1}}
\newcommand{\Webservice}[1]{\textsf{#1}}

\newcommand{\Refactoring}[1]{\Fachbegriff{#1}}
\newcommand{\CodeSmell}[1]{\Fachbegriff{#1}}
\newcommand{\Metrik}[1]{\Fachbegriff{#1}}
\newcommand{\DesignPattern}[1]{\Fachbegriff{#1}}
 % eigene projektspezifische Befehle, z.B. Abkürzungen usw.

\begin{document}

% \phantomsection
% \thispagestyle{empty}
% \pdfbookmark[1]{Eidesstattliche Erklärung}{ihkdeckblatt}
% \includegraphicsKeepAspectRatio{DeckblattIHK}{1}
% \cleardoublepage

\phantomsection
\thispagestyle{plain}
\pdfbookmark[1]{Deckblatt}{deckblatt}
% !TEX root = Projektdokumentation.tex
\begin{titlepage}

\begin{center}
\includegraphics[scale=0.25]{LogoIHK.pdf}\\[1ex]
\Large{Abschlussprüfung \pruefungstermin}\\[3ex]

\Large{\ausbildungsberuf}\\
\LARGE{\betreff}\\[4ex]

\huge{\textbf{\titel}}\\[1.5ex]
\Large{\textbf{\untertitel}}\\[4ex]

\normalsize
Abgabetermin: \abgabeOrt, den \abgabeTermin\\[3em]
\textbf{Prüfungsbewerber:}\\
\autorName\\
\autorAnschrift\\
\autorOrt\\
\autorNummer\\[5ex]

\includegraphics[scale=0.4]{\betriebLogo}\\[2ex]
\textbf{Ausbildungsbetrieb:}\\
\betriebName\\
\betriebAnschrift\\
\betriebOrt\\[5em]
\end{center}

% \small
% \noindent\\[1em]
% Dieses Werk einschließlich seiner Teile ist \textbf{urheberrechtlich geschützt}.
% Jede Verwertung außerhalb der engen Grenzen des Urheberrechtgesetzes ist ohne
% Zustimmung des Autors unzulässig und strafbar. Das gilt insbesondere für
% Vervielfältigungen, Übersetzungen, Mikroverfilmungen sowie die Einspeicherung
% und Verarbeitung in elektronischen Systemen.

\end{titlepage}

\cleardoublepage

% Preface --------------------------------------------------------------------
\phantomsection
\pagenumbering{Roman}
\pdfbookmark[1]{Inhaltsverzeichnis}{inhalt}
\tableofcontents
\cleardoublepage

\phantomsection
\listoffigures

\phantomsection
\listoftables

\phantomsection
% \addcontentsline{toc}{section}{Quellcodeverzeichnis}
\lstlistoflistings

\newcommand{\abkvz}{Abkürzungsverzeichnis}
\renewcommand{\nomname}{\abkvz}
\section*{\abkvz}
\markboth{\abkvz}{\abkvz}
\addcontentsline{toc}{section}{\abkvz}
% !TEX root = Projektdokumentation.tex

% Es werden nur die Abkürzungen aufgelistet, die mit \ac definiert und auch benutzt wurden.
%
% \acro{VERSIS}{Versicherungsinformationssystem\acroextra{ (Bestandsführungssystem)}}
% Ergibt in der Liste: VERSIS Versicherungsinformationssystem (Bestandsführungssystem)
% Im Text aber: \ac{VERSIS} -> Versicherungsinformationssystem (VERSIS)

% Hinweis: allgemein bekannte Abkürzungen wie z.B. bzw. u.a. müssen nicht ins Abkürzungsverzeichnis aufgenommen werden
% Hinweis: allgemein bekannte IT-Begriffe wie Datenbank oder Programmiersprache müssen nicht erläutert werden,
%          aber ggfs. Fachbegriffe aus der Domäne des Prüflings (z.B. Versicherung)

% Die Option (in den eckigen Klammern) enthält das längste Label oder
% einen Platzhalter der die Breite der linken Spalte bestimmt.
\begin{acronym}[WWWWW]
	\acro{API}{Application Programming Interface}
	\acro{CSS}{Cascading Style Sheets}
	\acro{HTML}{Hypertext Markup Language}\acused{HTML}
	\acro{IDE}{Integrated Development Environment}
	\acro{MVC}[MVC]{Model View Controller}
	\acro{SQL}{Structured Query Language}
	\acro{STI}{Single Table Inheritance}
	\acro{RoR}{Ruby on Rails}
	\acro{TDD}{Test Driven Development}
	\acro{LDAP}{Lightweight Directory Access Protocol}
	\acro{JSON}{JavaScript Object Notation}
\end{acronym}

\clearpage

% Inhalt ---------------------------------------------------------------------
\pagenumbering{arabic}
% !TEX root = Projektdokumentation.tex
% !TEX root = ../Projektdokumentation.tex
\section{Einleitung}
\label{sec:Einleitung}

\subsection{Projektumfeld}
\label{sec:Projektumfeld}
Die SUSE Linux GmbH wurde 1992 gegründet und ist seit 2014 eine Tochtergesellschaft der Micro Focus
AG. Weltweit beschäftigt das Unternehmen etwa 750 Mitarbeiter, welche auf die Standorte Nürnberg,
Prag, Peking und Provo (USA) verteilt sind. Der Hauptsitz der Firma befindet sich in Nürnberg,
ebenso wie der Hauptteil der Softwareentwicklung.

Das Kerngeschäft der SUSE Linux GmbH umfasst die Entwicklung einer Linux-basierten Distribution.
Für Privatkunden werden hierzu die Distributionen der openSUSE-Familie, die größtenteils von der
Community entwickelt werden, bereitgestellt, für Geschäftskunden die Produkte der SUSE Linux
Enterprise-Familie.

Die Entwicklung erfolgt nach dem Open-Source-Prinzip.
Durchgeführt wurde das Projekt im Team SUSE IT, die neben der Infrastruktur auch die
Entwicklungsumgebung, die sogenannten Engineering Services, bereitstellen.

\subsection{Projektbeschreibung}
\label{sec:Projektbeschreibung}
Zur Entwicklung der Open-Source-Software setzt die SUSE Linux GmbH in Teilen der Entwicklung
auf externe Tools und Services. Dies wird besonders im Bereich der Codeentwicklung und -pflege
deutlich, da neben den firmeneigenen Entwicklern auch Partner sowie Open-Source-Community
mitarbeiten. Hierbei kann es zu der Diskrepanz kommen, dass Entwickler neben ihrem Firmenaccount
auch private Accounts zur Entwicklung benutzen, wodurch eine genaue Zuordnung der Accounts und
ein Überblick über Zugriffsberechtigungen nahezu unmöglich wird.

Zur Lösung dieser Problematik soll eine Applikation entwickelt werden, welche es ermöglicht
eine Zuordnung der firmeninternen Accounts auf mögliche Accounts externer Services durchzuführen.
Das Ergebnis soll mittels einer API abrufbar sein.
Dabei soll die Applikation in ihrer Grundfunktion dem Gedanken der Self-Service-Technologies folgen,
damit Latenzen bei Änderungen der Accountinformationen minimiert und administratives Personal
entlastet wird. Um dies möglichst plattformunabhängig zu gestalten, erfolgt die Realisierung als
Webapplikation.

Da externe Services eine Zulassung durch den Betriebsrat und die IT-Sicherheit benötigen, beschränkt
sich die erste Version darauf, die Accountinformationen der Applikationen Github
(verzeichnisorganisierte Versionsverwaltung von Quellcode) und Trello (Tool zur Steuerung von
agilen Software-Projekten) einzubinden.

\subsection{Ist-Analyse}
\label{sec:Ist-Analyse}
Bei den verwendeten externen Services Github und Trello können Unternehmen zur Abgrenzung und
Verwaltung ihrer firmeninternen Ressourcen und Informationen sogenannte Organisationen anlegen, zu
denen Accounts von Mitarbeitern hinzugefügt werden und ihnen besondere Zugriffsrechte einzuräumen.
So haben beispielsweise nur Mitarbeiter Schreibrechte auf den Quellcodeverzeichnissen bei Github und
Mitglieder der Open-Source-Community sowie Partner können in der Regel nur Vorschläge, sogenannte
Pull Requests, bei dem jeweiligen Projekt einreichen.

Der Abgleich der internen Mitarbeiter- mit den Userlisten innerhalb der angelegten Firmenorganisationen
erfolgt bisher manuell. Dies bedeutet, dass beispielsweise für jeden neuen
Mitarbeiter persönlich oder schriftlich angefragt werden muss welches User-Alias er bei den
jeweiligen Services führt, um dessen Account dann zur Organisation hinzuzufügen und ihm
Zugriff zu ermöglichen. Ist dieser Aufwand bereits groß, erhöht er sich bei Mitarbeiteraustritten,
wenn es gegebenenfalls auf Grund fehlender Kontaktmöglichkeit es sehr mühsam wird eine Zuordnung
nachzuvollziehen.

\subsection{Soll-Konzept}
\label{sec:Soll-Konzept}

\subsection{Projektschnittstellen}
\label{sec:Projektschnittstellen}
Die Applikation nutzt mehrere APIs um z.B. die Betriebszugehörigkeit oder verschiedene Informationen
aus den Datenbanken der externen Services abzufragen.

Mitarbeiterspezifische Informationen, wie Name, Standort oder Vorgesetzter, werden aus dem
firmeneigenen eDirectory abgefragt.
Zur Abfrage der Organisationszugehörigkeit werden die Schnittstellen der genannten Tools Github und Trello
verwendet.

Die Schnittstelle der Applikation, welche die gesammelten Informationen über den Mitarbeiter
bereitstellt, wird in dem Administrationstool etsync konsumiert, welches für die Administratoren
der externen Organisationen bei Github und Trello entwickelt wurde.

% !TEX root = ../Projektdokumentation.tex
\section{Projektplanung}
\label{sec:Projektplanung}


\subsection{Projektphasen}
\label{sec:Projektphasen}
Für das Projekt wurde folgende Zeitplanung aufgestellt:

\tabelle{Zeitplanung}{tab:Zeitplanung}{ZeitplanungKomplett}

\subsection{Ressourcenplanung}
\label{sec:Ressourcenplanung}
\subsubsection{Personalplanung}
\label{sec:Personalplanung}
Das für die Entwicklung benötigte Personal besteht aus einem Auszubildenden, der dafür von seiner
regulären Mitarbeit im derzeitigen Team freigestellt wird.
\subsubsection{Kostenplanung}
\label{sec:Kostenplanung}
Die Personalkosten belaufen sich auf ein Auszubildendengehalt, wodurch lediglich Stromkosten für die
zur Projekterstellung eingesetzte Hardware anfallen. Durch die ausschließliche Nutzung einer
Entwicklungsumgebung und Werkzeugen aus dem Open-Source-Bereich entstehen keine zusätzlichen Kosten
wie Software-Lizenzen.

\subsection{Entwicklungsprozess}
\label{sec:Entwicklungsprozess}
Der Entwicklungsprozess erfolgt \textit{test-driven}, was bedeutet dass vor dem Schreiben des eigentichen
Quellcodes einer Funktionalität zunächst der dazugehörige Test entwickelt wird. Diese Vorgehensweise
ermöglicht einen abstrakteren und zielorientierteren Blick auf die eigentlichen geforderten
Funktionen der Applikation, in dem als erstes eine Überlegung über das Ergebnis einer Funktion
und daraufhin die eigentliche Implementierung des jeweiligen Algorithmus erfolgt, welcher zu dem
gewünschten Ergebnis führt.

\subsection{Architekturdesign}
\label{sec:Architekturdesign}
Zur Umsetzung der geforderten Modularität der Applikation bietet sich das sogenannte
\textit{Model-View-Controller-Schema} (\acs{MVC}) an. Dieses ist ein gängiges Schema beim Aufbau von
Webapplikationen und besteht aus drei Teilen:
\begin{itemize}
	\item das \textit{Model}, welches grob die einzelnen Tabellen einer Datenbank repräsentiert. Es enthält in
	der Regel zusätzliche Logik und Regeln, die zur Verwaltung und Zusammensetzung von Attributen
	eines Datenbankobjektes notwendig sind.
	\item der \textit{View}, welcher eine Ausgabe der gewünschten Daten bereitstellt. Im Kontext einer
	Webapplikation ist dies die eigentlich dargestellte Webseite in \textit{HTML/CSS}.
	\item der \textit{Controller}, welcher die Schnittstelle zwischen Model und View darstellt.
	Im Controller werden Eingaben verarbeitet und Befehle an das Model weitergegeben, welche
	die angeforderten Informationen an den Controller zurückgibt, der diese wiederum an den
	entsprechenden View verteilt.
\end{itemize}
Diese Aufteilung ermöglicht ein strukturierte und übersichtliche Entwicklung, da Funktionslogik,
Datenbankoperationen und Ausgabelogik klar getrennt und jeweils namentlich zugehörig gekennzeichnet
sind.

\abb{MVC-Schema}{fig:MVC-Schema}{MVC-Model.pdf}{0.5}

\subsubsection{Wahl der Programmiersprache}
\label{sec:Wahl der Programmiersprache}
Als Programmiersprache wurden \textit{Ruby} und das zugehörige Web-Framework \textit{Ruby on Rails}
gewählt.

Ruby basiert auf den Programmiersprachen Perl, Smalltalk, Eiffel, Ada und Lisp und legt Wert auf
die Balance von \textit{funktionaler} und \textit{imperativer Programmierung}. So verzichtet Ruby
beispielsweise auf Klammern zur Kennzeichnung von Codeblöcken und behandelt jeden Bestandteil des
Codes, wie Variablen,als ein \textit{Objekt}, welches eigene Methoden besitzt. Das macht die
Verarbeitung und Manipulation von Informationen sehr intuitiv und unterstreicht das Prinzip Rubys,
den Programmierer nicht durch Restriktionen und sprachlichen Eigenheiten bei seiner eigentlichen
Arbeit zu behindern.

Ruby und sein Framework Ruby on Rails werden quelloffen verwaltet und entwickelt. Neben den
existierenden Programmbibliotheken existieren zusätzlich von der Community bereitgestellte Module,
sogenannte \textit{Gems}, welche wie kleine zusätzliche Bibliotheken behandelt werden können. Die
Einbindung dieser Gems in ein bestehendes Projekt vermeidet Duplikation und erspart viel
Entwicklungszeit.

\subsection{Benutzeroberfläche}
\label{sec:Benutzeroberfläche}
Da ohne die Mitwirkung eines darauf spezialisierten \textit{UI/UX}-Designers möglichst zeiteffizient
eine für Nutzer verständliche Oberfläche realisiert werden musste, fiel die Entscheidung auf
\textit{Twitter Bootstrap}.
Das CSS-Framework Twitter Bootstrap wird häufig zur Gestaltung von Webseiten und
-applikationen verwendet, da es ein über eine umfangreiche Bibliothek an Gestaltungsvorlagen
verfügt und sehr leicht in ein Projekt zu integrieren ist. Auch unerfahrenen Anwendern gelingt damit
mühelos ein ansprechende und responsive Nutzeroberfläche.

\subsection{Rechteverteilung}
\label{sec:Rechteverteilung}
Die Applikation soll von zwei verschiedenen Arten von Usern verwendet werden: dem regulären
Mitarbeiter, der einzig Zugriff auf seine verwalteten Informationen zu externen Services hat, und
einem Administrator, der neben der Bearbeitung seines eigenen Datensatzes auch jene aller in der
Datenbank vorhandenen Mitarbeiter einsehen, abfragen und löschen kann.

Diese Rechteverteilung kann über ein \textit{Flag} gelöst werden, was die einfachste Methode ist.
Sie ist allerdings aus sicherheitsrelevanter Sicht sehr anfällig und leicht zu manipulieren,
weswegen entschieden wurde, die Rechteverteilung auf Datenbankebene zu realisieren.

\subsection{Datenmodell}
\label{sec:Datenmodell}
Das Projekt verfügt über zwei Informationsträger, die in der Datenbank abgebildet werden
sollen: einen \textit{User} und ein zugehöriges \textit{Tool}.

\abb{Entity-Relationship-Modell}{fig:ER-Modell}{ER-Diagramm.pdf}{0.4}

Da für die Applikation zusätzliche Spezialisierungen, wie eine Administrator-Rolle und verschiedene
externe Services, notwendig sind, wurden diese Beziehungen mit Hilfe der sogenannten
\textit{Single Table Inheritance} (\acs{STI}) realisiert.

\abb{Beispielhafte Darstellung der STI}{fig:STI-Beispiel}{STI-Figure.pdf}{0.5}

Diese spezielle Vererbungsstrategie lässt sich mit Hilfe der in Rails eingebauten
\textit{Active Record}-Engine durch das einfache Hinzufügen eines Typ-Attributs an das Model
umsetzen.

\subsection{Schnittstellen}
\label{sec:Schnittstellen}
Die Hauptaufgabe der Applikation ist die Bereitstellung der gesammelten Informationen zu den
Mitarbeitern und deren eingetragenen Informationen zu ihren genutzten Services. Dazu soll eine
eigene Schnittstelle innerhalb der Applikation entwickelt werden, welche die Informationen im
\textit{JSON}-Format ausgibt. Diese Schnittstelle kann nur authentifiziert angesprochen werden,
\dahe der User muss eine Administratorrolle besitzen, um Zugriff zu erhalten.

\subsection{Paketierung}
\label{sec:Paketierung}
Gemäß der Richtline des Teams, in dem die Applikation entwickelt wurde, wird jene als \textit{\acs{RPM}}
paketiert. Dies ermöglicht eine einfache und sofort funktionsfähige Installation, da alle
notwendigen Programmabhängigkeiten im Paket definiert und während des Installationsprozesses
abgerufen und mitinstalliert werden. Dadurch ist beispielsweise auch bei Bedarf eine problemlose
Migration auf einen anderen Server möglich.

% !TEX root = ../Projektdokumentation.tex
\lstset{
  backgroundcolor=\color{white}
}

\section{Implementierungsphase}
\label{sec:Implementierungsphase}

\subsection{Einleitung}
\label{sec:Einleitung}
Nach der Planungsphase wird im direkten Anschluss die Implementierungsphase eingeleitet. Diese
umfasst die folgenden Schritte, welche in diesem Kapitel einzeln näher erläutert werden:
\begin{enumerate}
	\item Implementierung des Backends
	\item Implementierung der Benutzeroberfläche
	\item Anbindung externer Schnittstellen
	\item Implementierung der internen Schnittstelle
	\item Durchführung einer Testphase
	\item Paketierung der Applikation
	\item Deployment auf Produktivmaschine
\end{enumerate}

\subsection{Implementierung des Backends}
\label{sec:Implementierung des Backends}
Als erster Schritt erfolgt die initiale Erstellung eines Ruby on Rails-Projektes.
Nach Installation der zugehörigen Entwicklungsumgebung geschieht dies einfach mit Ausführung des
Befehls \Code{rails new} im Arbeitsverzeichnis. Mit Hilfe dieses Befehls wird eine
standardisierte Projektmappe erzeugt.

RoR besitzt viele integrierte Kommandozeilenbefehle, mit denen es
möglich ist, die Ausführung von
Entwicklungsschritten zu automatisieren. Beispielsweise kann das Anlegen
von relevanten Dateien innerhalb der entsprechenden Verzeichnisstruktur
direkt über diese Befehle erledigt werden, was zu einer Beschleunigung
der Entwicklungsprozesse führt.

Das Gemfile bildet die zentrale Auflistung der von der Applikation genutzten Gems. Da die
Entwicklung der Applikation im Sinne des \acs{TDD} (\Vgl \ref{sec:Entwicklungsprozess}) erfolgt,
wird hier zunächst die Testsuite \anf{rspec} hinzugefügt und die ersten Tests für die Models
(\Vgl \ref{sec:Architekturdesign}) geschrieben.

\subsubsection{Erstellen der Models}
\label{sec:Erstellen der Models}
Erst nach Schreiben der Tests erfolgte die Erstellung der eigentlichen Model-Dateien samt jeweils
zugehörigen Attributen. Dies geschieht mit Hilfe sogenannter Migrationen,
Datenbankoperationen geschrieben in Ruby. Zum Hinzufügen der so definierten Tabellen und Attribute
zur Datenbank wird der Befehl \Code{rake db:migrate} ausgeführt, welcher die Ruby-Syntax in
\acs{SQL} übersetzt.

Um die Integrität der Datenbanktabellen zu erhalten und \ggfs Fehleingaben der Anwender abzufangen,
werden in den Models Validierungen implementiert. Diese werden bei jeder Datenbankoperation
ausgeführt und sollen verhindern, dass fehlerhafte Daten in der Datenbank abgespeichert werden.

\snip{externaltool.rb}{1}{6}

Im Anschluss wird eine weitere Migration ausgeführt, welche ein Typ-Attribut zu den
beiden Models hinzufügt. Hiermit wird die in der
Projektplanung erfasste Vererbungsstrategie \acs{STI} realisiert (\Vgl
\ref{sec:Datenmodell}).

\subsubsection{Erstellen der Controller}
\label{sec:Erstellen der Controller}
Entsprechend der Konvention wird für jedes Model mindestens ein
Controller erstellt, jedoch benötigt nicht jeder Controller ein Model.
Allerdings ist in jedem Controller mindestens eine Methode, die als
\anf{Action} bezeichnet wird, zwingend erforderlich. Wird beispielsweise
durch einen Benutzer die Webseite zum Hinzufügen eines neuen Alias
aufgerufen, wird diese Anfrage an die \Code{new}-Action des
entsprechenden Controllers weitergeleitet.
(\Vgl \ref{sec:Architekturdesign})

Den Einstiegspunkt der Applikation bildet der \anf{Application Controller}. In diesem werden
verschiedene Strategien implementiert, welche den Betrieb der gesamten Applikation betreffen,
beispielsweise das initiale Aufrufen der Nutzerauthentifizierung. Alle anderen Controller erben von diesem
Application Controller und verfügen damit bei Aufruf über alle dort niedergelegten Methoden.

Zunächst wird der \anf{Externaltools Controller} erstellt. Alle relevanten Operationen
wie Anzeige, Hinzufügen und Löschen von Daten des Tool-Models werden von diesem Controller
verarbeitet. Dazu fragt der Controller Informationen zu dem eingeloggten Nutzer aus der User-Tabelle
ab. Dies wird über das Objekt \Code{current user} realisiert. Das Objekt stammt aus dem
eingesetzten Authentifizierungs-Gem \anf{Devise} (\Vgl \ref{sec:Anbindung externer Schnittstellen}).

\snip{externaltools.controller.rb}{18}{33}

Zur Abgrenzung des Administratorbereichs der Applikation wird ein Administratorverzeichnis angelegt.
Um die Controller der dort zu implementierenden Funktionen vor einem unberechtigten Zugriff zu schützen, wird ein
\anf{Admin Controller} erstellt. Dessen einzige Methode prüft, ob ein Administrator eingeloggt ist.
Ist dies nicht der Fall, erfolgt eine Weiterleitung zur Index-Seite.

Da nur die Administratoren die Möglichkeit besitzen sollen, Nutzer aus der Datenbank zu entfernen,
befindet sich der \anf{Users Controller} in diesem Unterverzeichnis. Als Kernstück dieses
Controllers wird die \Code{list user}-Action implementiert, welche die gesammelten Informationen
zu allen Mitarbeitern aus der Datenbank abfragt und die API bereitstellt.

\snip{users.controller.rb}{26}{29}

\subsection{Implementierung der Benutzeroberfläche}
\label{sec:Implementierung der Benutzeroberfläche}
Nach der Fertigstellung des grundlegenden Backends erfolgt die Programmierung der Benutzeroberfläche.
Auch für diese werden Tests geschrieben, die sogenannten \textit{Feature Tests}.

\snip{employee.signin.spec.rb}{11}{15}

Wie in der Projektplanung (\Vgl \ref{sec:Benutzeroberfläche}) niedergelegt, wird für die Implementierung
das Framework Twitter Bootstrap verwendet.

Die Nutzeroberfläche lässt sich zeitsparend mit den in Bootstrap vorhandenen Vorlagen erstellen. Die Farbgebung der Elemente wird gemäß der
Firmenrichtlinien gestaltet.

\abb{Nutzeroberfläche}{fig:Nutzeroberfläche}{externaltoolsscreen.pdf}{0.7}

Als Strategie zur Verbesserung der Performance der Applikation werden sogenannte
\textit{Partial Views} eingesetzt. Bei dieser Herangehensweise werden bei Bedarf nur die sich
verändernden Elemente auf einer Webseite neu geladen, statt die komplette Seite neu abzurufen.
Dies verringert die benötigte Antwortzeit des Webservers erheblich.

Für Zugang zu den Administratoransichten wird die reguläre Mitarbeiteransicht um weitere Tab-Elemente
erweitert. Diese sind nur bei Login als Administrator sichtbar.

\subsection{Anbindung externer Schnittstellen}
\label{sec:Anbindung externer Schnittstellen}
Als Authentifizierungsmethode wird eine Authentifizierung über \acs{LDAP} gewählt. Damit
kann das firmeneigene eDirectory zum Abgleich genutzt und gleichzeitig Informationen über
den Mitarbeiter abgerufen werden.

Das Gem \anf{Devise} bietet eine umfassende Sammlung an Konfigurationsmöglichkeiten und wird
daher als Lösung eingesetzt. Beim erstmaligen Einloggen eines Mitarbeiters werden gleichzeitig
dessen Name, Mitarbeiterkürzel und Informationen zu seinem Standort sowie die interne
e-Mail-Adresse in der Datenbank abgespeichert. Informationen zum jeweiligen Vorgesetzten des
Nutzers werden einzeln bei Bedarf abgefragt und angezeigt.

Der Beschäftigungsstatus eines Mitarbeiters soll möglichst aktuell als Information vorliegen. Da diese
Information allein in der Administratoransicht \bzw über die Schnittstelle abgefragt wird,
würde eine gesammelte Abfrage des Status für mehrere hundert Mitarbeiter gleichzeitig den LDAP-Server
überlasten. Aus diesem Grund wird ein Hintergrundprozess implementiert. RoR bietet dafür den
sogenannten \textit{Active Job}-Mechanismus an. Dort wird ein Prozess als Klasse implementiert,
welcher einmal pro Woche für alle eingetragenen Mitarbeiter einzeln den Mitarbeiterstatus anfragt
und als Flag in der Datenbank aktualisiert.

\snip{employee.status.job.rb}{1}{13}

Als Kernstück wird die Klasse \Code{LDAPSearch} implementiert, welche für alle zusätzlichen
LDAP-Abfragen außerhalb des User-Models verwendet wird.

\snip{ldap.search.rb}{4}{14}

Gewünscht war ebenfalls die Möglichkeit, dass Mitarbeiter \ggfs in Eigeninitiative eine
Beitrittsanfrage an die Administratoren der Organisationen stellen können. Zu diesem Zweck wird der
Status der Organisationszugehörigkeit ebenfalls durch einen ActiveJob bei der API von Github
angefragt, welcher beim Hinzufügen eines Github-Nutzernamens ausgelöst wird.
Bei nicht vorhandener Zugehörigkeit wird ein Link bereitgestellt, welcher eine fertige e-Mail-Vorlage
an die Administratoren generiert.

Aus Zeitmangel konnte eine ähnliche Strategie für den externen Service Trello während des
Projektverlaufs nicht mehr implementiert werden. Da dies aber nicht als kritisch bewertet wurde,
behinderte dies nicht den Abschluss des Projektes.

\subsection{Interne Schnittstelle}
\label{sec:Interne Schnittstelle}
Die Schnittstelle wird für die Verarbeitung durch das Kommandozeilenwerkzeug \anf{etsync} entwickelt.
Anhand des bestehenden Codes in \anf{etsync} wird die Ausgabe der Schnittstelle modelliert.
Wie bereits in der Planungsphase erwähnt, erfolgt die Realisierung in \acs{JSON},
welches ein gebräuchliches Datenformat zum Austausch von Informationen zwischen Anwendungen ist.

Um die gewünschten Informationen korrekt und einfach zu modellieren, wird das Gem \textit{rabl}
eingesetzt. Mit diesem Gem kann die JSON-Notation einfach als Datenbankabfragen in Ruby
geschrieben werden und vermeidet so Fehler bei der sonst manuellen Erstellung einer JSON-Vorlage,
welche viele Klammern und Satzzeichen verwendet. Auch können spätere Anpassungen so leichter
durchgeführt werden da eine genaue Kenntnis der JSON-Syntax überflüssig ist.

Die zugehörige Controller-Action befindet sich im Administratorverzeichnis, weshalb auf die Schnittstelle
nur von Administratoren authentifiziert zugegriffen werden kann (\Vgl \ref{sec:Erstellen der Controller}).

\subsection{Testphase}
\label{sec:Testphase}
Nach Implementierung des Projektes konnte eine finale Testphase zeitlich gering gehalten werden.
Durch den Einsatz von \acs{TDD} wurde die Applikation während der Entwicklung bereits laufend
getestet und eventuelle Fehler wurden behoben. Neben der Durchführung von \textit{Usability-Tests}
der Nutzeroberfläche mit Mitarbeitern und ihren Daten fand abschließend ein kompletter
Durchlauf der Testsammlung statt.

\subsection{Paketierung der Applikation}
\label{sec:Paketierung der Applikation}
Wie bereits in der Projektplanung dargelegt erfolgt nach Entwicklung und finaler Testphase der
Applikation die Paketierung als RPM. Die Kontrollstruktur eines RPM-Pakets ist das sogenannte
Specfile. In diesem werden alle relevanten Informationen zu dem Softwarepaket,
wie beispielsweise Versionsnummer oder Abhängigkeiten auf andere Pakete, niedergelegt.

Um das Paket immer auf dem neuesten Stand zu halten, wird über ein Skript bei einer
Codeaktualisierung des Projekts automatisch die Ausführung der Testsammlung ausgelöst.
Sind diese erfolgreich mit den neuesten Änderungen ausgeführt worden, wird die Paketierung
automatisiert angestoßen. Dies verhindert zusätzlich, dass fehlerhafte
oder ungetestete Funktionalität in das endgültige Paket gelangt.

\subsection{Deployment}
\label{sec:Deployment}
Die Bereitstellung der Applikation geschieht in Kollaboration mit dem Team, in welchem das
Projekt durchgeführt wurde. Auf einer Produktivmaschine, die bereits als Webserver für
andere Applikationen fungiert, wird eine weitere Konfiguration angelegt und das Software-Paket der
Applikation installiert. Nach Neustart der zugehörigen Services und Prozesse ist die
Webapplikation intern über den Browser zu erreichen.
\pagebreak

% !TEX root = ../Projektdokumentation.tex
\section{Projektabschluss}
\label{sec:Projektabschluss}

\subsection{Soll-/Ist-Vergleich}
\label{sec:SollIstVergleich}

\paragraph{Beispiel (verkürzt)}
Wie in Tabelle~\ref{tab:Vergleich} zu erkennen ist, konnte die Zeitplanung bis auf wenige Ausnahmen eingehalten werden.
\tabelle{Soll-/Ist-Vergleich}{tab:Vergleich}{Zeitnachher.tex}

\subsection{Projektabnahme}
\label{sec:Projektabnahme}
Die Projektabnahme erfolgte mit einer kurzen Präsentation und Demonstration der Applikation für den
Projektbetreuer sowie die in Frage kommenden Administratoren der Organisationen.

\subsection{Firmenweite Einführung}
\label{Firmenweite Einführung}
Nach Abnahme des Projekts wurde zu einem festgelegten Stichtag eine Ankündigungsmail auf den
internen Mailinglisten verschickt, welche kurz den Nutzen der Applikation erklärt und alle
Mitarbeiter zur Eintragung ihrer Daten aufruft. Desweiteren wurde das Anlegen des Datensatzes für
jeden neuen Mitarbeiter verpflichtend als Richtlinie eingeführt.

% !TEX root = ../Projektdokumentation.tex
\section{Fazit}
\label{sec:Fazit}
In der vorgegebenen Zeit konnten alle definierten Projektziele umgesetzt werden. Bis auf
Verzögerungen bei der Anbindung externer Schnittstellen sind alle relevanten Funktionen der
Webapplikation implementiert worden.
Nach der firmenweiten Einführung fanden bereits mehrere Überprüfungen der Mitgliederlisten in den
externen Organisationen statt und der Prozess konnte, wie bereits dargestellt, erheblich beschleunigt werden.
Statt der vorherigen
zeitintensiven Informationsbeschaffung genügt nun die Ausführung eines Befehls auf der Kommandozeile,
welcher die gesammelten Daten aus der Webapplikation abruft. Das Projekt kann somit also als
erfolgreich gewertet werden.

\subsection{Ausblick}
\label{sec:Ausblick}
Die Webapplikation wird nach Abschluss weiterhin von mir als Hauptentwickler gepflegt und laufend
gewartet.

Folgende zusätzlichen Funktionen sind geplant oder vorgeschlagen worden:
\begin{itemize}
  \item Einbindung in die Authentifizierungsinfrastruktur der Micro Focus PLC, ähnlich des
  Intranets
  \item Vergabe des Administratorstatus an Mitarbeiter innerhalb der Applikation
  \item Auslagerung aller administrativen Arbeiten an die Webapplikation (Abgleich der
  Mitarbeiterlisten, Hinzufügen und Entfernen von Mitarbeitern in den externen Organisationen)
  \item Anbindung an die Trello API
\end{itemize}



% % Anhang ---------------------------------------------------------------------
% \clearpage
% \appendix
% \pagenumbering{roman}
% \input{Anhang}

\end{document}
